
We set out to measure a precision hierarchy. The intuition seemed safe: the more formally rigorous the verifier sitting in an LLM's repair loop, the more precisely it should describe what went wrong, and the better the model should repair its own code. We ran four formal feedback categories over the same failing solutions with the same backbone and the same repair budget, ranked them by how much feedback text they produced, and correlated that ranking against repair success. The Spearman correlation came back at exactly -1.0. The most informative verifier fixed the fewest programs.

The numbers behind that inversion are stark. Pyright hands GPT-4o-mini an average of 24,358 characters of structured JSON diagnostics per problem and repairs 4.76\% of failures on the first iteration. A 202-character Python traceback repairs 5.56\%. Then wall-clock cost enters, and the picture inverts a second time: because Pyright runs in roughly 46ms against execution monitoring's 800ms, static analysis returns an estimated 6.64 points of pass@1 per second of verifier time while execution monitoring returns 0.41. The verifier that repairs the most code per attempt is not the verifier that repairs the most code per second, and neither is the one the specificity intuition would have predicted.

This matters because the intuition is load-bearing. Systems that wrap an LLM in a Generate-Verify-Repair loop have to choose what goes in the verify stage, and that choice is currently made on an implicit ranking --- SMT counterexamples above static type errors above execution traces --- that, as far as we can determine, has never been measured end to end.

At the surface, the problem is familiar and largely solved: LLMs write incorrect code, and feeding verifier output back into the prompt recovers some of it. Self-Repair \citep{Olausson2023SelfRepair}, Reflexion \citep{Shinn2023Reflexion}, and CodeT \citep{Chen2022CodeT} all establish gains in the 5--15\% range on HumanEval and MBPP, and the mechanism is no longer in doubt.

The deeper problem is what those results cannot tell us. Each studies a single feedback category, on its own backbone, against its own baseline, and none reports what the feedback cost to produce. Feedback specificity and feedback cost are separate axes, and the field has been reasoning about the first while paying for the second. Comparing categories demands that the backbone, the benchmark, the prompt template, and the repair budget all be held fixed while only the verifier changes --- an experimental design that no single-method paper had reason to run, because each had a method to advocate rather than a taxonomy to characterize.

That leaves a concrete gap: there is no overhead-normalized, backbone-controlled comparison of formal feedback categories for LLM code repair. Without one, verifier selection rests on an assumption about formal rigor, and correctness-only reporting makes an expensive verifier look free. Our results suggest the assumption is not merely weakly supported but reversed at the model tier most practitioners deploy.

Our explanation is a single shift in what bounds repair. LLM repair is limited by the model's ability to extract one actionable signal from the feedback, not by how much information the feedback contains. A traceback naming a line and an exception type is actionable on sight. Pyright's diagnostic tree, even after truncation to 4,000 characters to fit the prompt, buries the decisive diagnostic among unused-import notices and missing-annotation warnings. Information content rises; extractable signal falls. And because a tool's verbosity has nothing to do with its runtime, the cheapest verifier can be the one with the best cost-normalized return --- which is what we observe.

Working from that insight, this paper contributes the following. We report the first controlled comparison of four formal feedback categories --- execution monitoring, static analysis, type checking, and SMT solving --- on a single backbone (GPT-4o-mini) over 421 HumanEval and MBPP problems, with wall-clock overhead measured alongside correctness. We characterize the failure population that any such loop has to repair, finding a mixed bug distribution in which logic errors dominate at 65.9\% and no single type exceeds the 80\% that would make the comparison degenerate. We establish that the categories produce genuinely different signals rather than paraphrases of each other (Kruskal-Wallis H=338.78, p=4.$01\times 10^{-73}, \epsilon ^{2}=0.88)$, which is the precondition that makes the comparison meaningful. We then measure the inverse relationship between feedback length and single-iteration repair utility that gives the paper its title, and normalize by overhead to show that static analysis dominates on correctness per second. Finally, we report that SMT-guided repair is not deployable at this model tier at all: GPT-4o-mini extracted usable Z3 constraints for 6.3\% of problems against the 40\% our design assumed, a negative result that marks a boundary for anyone planning an SMT stage.

Two of these findings emerged from experiments that failed their pre-registered gates. We had predicted a positive specificity-repair correlation and an efficiency win for execution monitoring; we got neither. We report them as the paper's central results because a refuted prediction that reverses a field-wide assumption is more useful than a confirmed one that restates it.

Section 2 positions this work against prior feedback-loop and formal-repair literature. Section 3 describes the experimental design and the controls that make the comparison valid. Section 4 states the research questions and setup. Section 5 presents results. Section 6 discusses mechanisms, confounds, and deployment recommendations.

